\section{Shared tensor contractions for the complex side map}
\label{sec:disjoint-tensor-compression}
This section changes only the disjoint-triple side circuit of
Section~\ref{sec:complex-compression}. The scalar field, intersection-two
circuit, centers, twelve signed operations, tensor stages, endpoint
corrections, and first/third-stage matching remain as there. It supplies
complex-network headroom, not a new multiplication exponent.

\subsection{An exact sum recurrence}
Let $D_{a,b}(n)$ send values on $a$-subsets to their sums over subsets
disjoint from each specified $b$-subset. Split the ground set into $l+r=n$.
For sources with left size $i$ and targets with left size $j$, the map is
$D_{i,j}(l)\otimes D_{a-i,b-j}(r)$. Applying its two factors in either order
uses at most
\[
 \min\left\{\binom r{a-i}c(l,i,j)+\binom lj c(r,a-i,b-j),\quad
 \binom li c(r,a-i,b-j)+\binom r{b-j}c(l,i,j)\right\}
\]
additions. Here $c$ is the already constructed smaller circuit's addition
bound. Sum these costs over feasible $i,j$, then add
$\binom lj\binom r{b-j}(p_j-1)$ for each $j$ with $p_j>0$ feasible source
splits. Try every split up to $n/2$, as well as the direct row-sum circuit.
Empty maps, copies, and the one-output total are the base cases.
All combined supports are disjoint: the tensor factors partition sources
by their two parts, and the final sum partitions them by $i$.
This proves the integer coefficient map inductively.

Intern equal formal supports, retaining the first decomposition, and prune
unused nodes. These operations preserve the map and do not increase the
number of additions. Each noninput retained node is still the sum of two
disjoint nonempty supports. At $h=26$ the unshared recurrence gives $74162$
additions; exact interning and pruning give $57772$ additions and $2600$
outputs, hence $60372$ reversible roles by the $c+q$ embedding.

\subsection{A phase enclosure for any disjoint-triple sum DAG}
For a retained node $z$, let $\mathcal S_z$ be its ancestor input triples
and $\mathcal T_z$ the output triples reachable from it. Every pair in
$\mathcal S_z\times\mathcal T_z$ is disjoint. Both families are nonempty.
Write $A_z=\bigcup\mathcal S_z$ and $B_z=\bigcup\mathcal T_z$; thus
$A_z\cap B_z=\varnothing$. In the forward middle mixer assign
\[
 U_z=\begin{cases}
 \langle t_S\rangle,&\mathcal S_z=\{S\},\\
 t_T^\perp,&|\mathcal S_z|>1,\ \mathcal T_z=\{T\},\\
 \mathbb F_2^{A_z},&|\mathcal S_z|,|\mathcal T_z|>1.
 \end{cases}
\]
Use the same rule with $\mathcal S,\mathcal T$ exchanged for the reverse
middle mixer. This is not obtained by assuming complements preserve
nonalternation.

Along a forward DAG edge, ancestor families grow and reachable-output
families shrink. Consequently the displayed spaces increase. Lines remain
the same until a second source appears; complements remain the same once
only one target remains. Coordinate spaces increase by inclusion. A line
can enter a coordinate enclosure or target complement, and a coordinate
enclosure can enter a target complement. There are no other distinct
internal transitions. The same proof applies in reverse with the families
exchanged. Initial physical source lines and final physical target
complements fit these inclusions, including nodes that are themselves
outputs and have other later uses.

Every space is nondegenerate under the standard binary dot product.
A nonzero coordinate residual has a coordinate unit. For a source-line
to coordinate transition, two distinct ancestor triples ensure a point in
$A_z\setminus S$, giving a norm-one residual vector. For a coordinate
space on $A$ entering $t_T^\perp$, the earlier node has another reachable
triple $T'\ne T$ disjoint from $A$. A point in $T'\setminus T$ gives a
coordinate unit orthogonal to both $A$ and $T$. A line-to-complement
transition has a coordinate unit outside $S\cup T$, since $h>6$.
Transitions from zero or to the full space likewise have coordinate units,
except the complement-to-full residual, which is the norm-one line $t_T$.
Thus every nonzero orthogonal residual is nonalternating. The retained
binary-space lemma constructs an orthonormal basis for each residual.
Tensoring and the three-stage frame schedule are consequently licensed
exactly as in Section~\ref{sec:complex-compression}.

The explicit checker traverses every physical role segment of both middle
orientations, including source, target, fresh, and retired roles. At small
sizes it also constructs every distinct residual basis independently of
the case proof. The scalar checker verifies every output coefficient over
the integers, and the signed dirty-workspace invocation is checked on
independent formal inputs in both directions.

\subsection{Complete finite accounting}
With $v=\binom h3$, retain $R_2=\binom h2(4(h-2)-6)$ intersection-two roles.
Let $R_0=c+v$ be the new disjoint roles. Then
\[
 W=2v^3+2v^2(R_0+R_2+h+1),\quad m=h^3,\quad
 L=3v^2h(h+1),\quad s=Wm-2v^3+2L.
\]
All side and central inputs are arbitrary and restored. The join uses the
same paired-point label permutation as before; its binary residual has a
norm-one coordinate tensor outside the two relevant triple supports.
It removes one auxiliary role and $m$ rank per join without further loss.
At $h=26$, $R_0+R_2=89622$ and $Wm-s=6678880000$; exact logarithm bounds
certify $1-\log_m(s/W)>3\cdot10^{-8}$.

If the mixer uses $M$ scalar additions, has $C$ data copies and $J$ weighted
output injections, the three-stage scalar step bound remains
$3v^2(4M+2C+4J+18v)$. The checker computes these numbers from the expanded
DAG. With $E=64(W+m+1)^3$, it verifies
$4\,\mathrm{steps}+8s+8W+8m<E$ and $2\le s<m^5$.
Thus the existing stopped-depth guard lemma applies with this enlarged
complex saving and the same residual exponent five. This check changes no
retained multiplication witness.
